\documentclass[11pt]{article}
\usepackage[margin=1.6cm]{geometry}
\usepackage[T2A,T1]{fontenc}
\usepackage[utf8]{inputenc}
\usepackage[russian]{babel}
\usepackage[hidelinks]{hyperref}
\usepackage{enumitem}
\usepackage{tabularx}

\setlist[itemize]{nosep, leftmargin=1.2cm, label=\textbullet, topsep=2pt, itemsep=2pt}

\pagenumbering{gobble}

\begin{document}

% --- HEADER ---
\begin{center}
    {\LARGE \textbf{Гурьянова Ксения}}\\[4pt]
    \large ML Engineer / Data Scientist
\end{center}

\begin{center}
    \href{mailto:keguryanova@gmail.com}{keguryanova@gmail.com} \,|\, 
    \href{https://github.com/KseniaGu}{github.com/KseniaGu}
\end{center}

\vspace{10pt}

% --- PROFILE ---
\noindent \textbf{Профиль:} Опыт работы более 4 лет в сфере разработки и внедрения решений для e-commerce и аналитических систем. Владею полным циклом ML-разработки (сбор — обработка данных — обучение — оптимизация — деплой). Основной фокус: обработка естественного языка (NLP). Имею опыт работы с компьютерным зрением (CV), навыки работы с СУБД, а также опыт разработки Graph RAG-системы в рамках практического проекта (\href{https://github.com/KseniaGu/remnote-graph-rag}{ссылка}).

\vspace{10pt}

% --- EXPERIENCE ---
\section*{Опыт работы}

\noindent
\begin{tabularx}{\textwidth}{@{}l X@{}}
\textbf{2021--2025} & \textbf{ML Engineer} \,|\, \href{https://scented.ai/}{scented.ai}
\end{tabularx}

\vspace{2pt}
\noindent \textit{AI-стартап, разрабатывающий решения для парфюмерной индустрии}

\vspace{4pt}
\noindent \textbf{Основные задачи:}
\begin{itemize}
  \item Разработала пайплайн сопоставления товаров (product matching) с использованием ML-подходов; работала с эмбеддингами для текстовых данных и изображений.
  \item Настроила поисковую систему для текстовых данных (BM25, text normalization).
  \item Разработала систему обнаружения клонов продуктов на основе отзывов (классический ML).
  \item Решала задачи ранжирования продуктов по релевантности в разных категориях.
  \item Проводила анализ пользовательского фидбэка для настройки алгоритмов отбора наиболее релевантного контента.
  \item Участвовала в разработке REST API-сервиса (FastAPI, Docker); отвечала за логирование.
  \item Реализовывала интеграцию и поддержку баз данных: PostgreSQL (логи, пользовательские данные), MongoDB (история обновлений системы), Redis (кеширование часто запрашиваемых данных).
\end{itemize}

\vspace{10pt}

\noindent
\begin{tabularx}{\textwidth}{@{}l X@{}}
\textbf{2023--2025} & \textbf{ML Education Specialist} \,|\, \href{https://school.razinkov.ai/}{school.razinkov.ai}
\end{tabularx}

\vspace{2pt}
\noindent \textit{AI онлайн-школа, посвященная подготовке специалистов в области машинного обучения}

\vspace{4pt}
\noindent \textbf{Основные задачи:}
\begin{itemize}
  \item Участвовала в разработке практических заданий по базовому и продвинутому глубокому обучению: классические алгоритмы, свёрточные нейронные сети (CNN, ResNet, U-Net), архитектура трансформер (Transformer, ViT).
  \item Участвовала в составлении рекомендаций по практикам MLOps: мониторинг (Neptune.ai/MLflow), вывод моделей в эксплуатацию (Docker, FastAPI, ONNX).
  \item Проводила эксперименты: классификация и регрессия (изображения, табличные данные), поиск похожих изображений (image retrieval), сегментация изображений (semantic segmentation), машинный перевод (machine translation).
  \item Проводила код-ревью.
\end{itemize}

\vspace{10pt}

% --- TECHNICAL SKILLS ---
\section*{Технические навыки}

\noindent \textbf{Языки программирования:} Python; SQL. Изучались: C++/C\#\\[4pt]
\textbf{Машинное обучение:}
\begin{itemize}
  \item \textbf{Фреймворки:} PyTorch, scikit-learn. Есть опыт с TensorFlow, Keras.
  \item \textbf{Библиотеки:} Transformers, Sentence Transformers, NumPy, Pandas, spaCy, NLTK.
  \item \textbf{LLM:} практический опыт применения LLM (GPT, Gemini, Mistral, Qwen) и их развёртывания через vLLM/Ollama, а также разработки Graph RAG‑системы (LangGraph, LlamaIndex, Cloud Run).
  \item \textbf{Исследовательский опыт:} архитектура Transformer (Encoder–Decoder, BERT, RoBERTa, Siamese BERT, ViT), CNN (VGG, ResNet, U‑Net), RNN (LSTM/ESIM), GAN, классический ML (логистическая/линейная регрессии, SVM, деревья решений и т.п.).
\end{itemize}
\vspace{4pt}
\noindent \textbf{Базы данных и поиск:} PostgreSQL, Redis, MongoDB, ClickHouse, Faiss, Meilisearch.\\[4pt]
\textbf{MLOps и развёртывание:} Docker, FastAPI, Git (GitHub (CI/CD с Actions), GitLab (CI/CD)), Neptune.ai, MLflow, ONNX, TorchScript, Google Cloud Platform (GCP).\\[4pt]

\vspace{10pt}

% --- EDUCATION ---
\section*{Образование}

\noindent
\begin{tabularx}{\textwidth}{@{}l X@{}}
\textbf{2019--2021} & \textbf{Высшее образование (магистратура)} \,|\, КФУ, Институт вычислительной математики и информационных технологий, направление «Фундаментальная информатика и информационные технологии», профиль «Машинное обучение и компьютерное зрение».
\end{tabularx}

\vspace{6pt}

\noindent
\begin{tabularx}{\textwidth}{@{}l X@{}}
\textbf{2014--2019} & \textbf{Высшее образование (бакалавриат)} \,|\, КФУ, Институт вычислительной математики и информационных технологий, направление «Фундаментальная информатика и информационные технологии», средний балл 5,0 (из 5,0).
\end{tabularx}

\vspace{6pt}

\noindent
\begin{tabularx}{\textwidth}{@{}l X@{}}
\textbf{2009--2014} & \textbf{Среднее (полное) общее образование} \,|\, МАОУ «Лицей №131», золотая медаль.
\end{tabularx}

\vspace{6pt}

\noindent \textbf{Сертификат о прохождении курса на Coursera:} \href{https://coursera.org/share/4f64abf86f701eb8df6e2b00cb6e4768}{«SQL for Data Science»}


\vspace{10pt}

% --- PROJECTS/THESIS ---
\section*{Выпускные работы}

\subsection* {Магистерская диссертация}
\textbf{Название:} «Исследование задачи определения корректности логического вывода на естественном языке с использованием машинного обучения»\\
\textbf{Описание:} реализована модель глубокого обучения для решения исследуемой задачи, разработан алгоритм дополнения данных с применением языковых моделей, собран набор данных для проведения экспериментов по обработке длинных текстов моделями глубокого обучения на основе архитектуры Transformer.\\
\textbf{Ссылки:} \href{https://github.com/KseniaGu/NLI-based-augmentation}{NLI-based data augmentation}, \href{https://github.com/KseniaGu/movie-screenplays/tree/main}{Movie screenplays}, \href{https://www.kaggle.com/datasets/gufukuro/movie-scripts-corpus}{Movie Scripts Corpus}

\vspace{8pt}


\subsection* {Выпускная квалификационная работа}
\textbf{Название:} «Применение методов машинного обучения при составлении аннотации текста»\\
\textbf{Описание:} разработан алгоритм автоматического реферирования текстов на русском языке, основанный на кластеризации векторных представлений предложений.

\end{document}